(1) 初态概率密度 \(|\psi(x,0)|^2=(2\pi\sigma_0^2)^{-1/2}e^{-x^2/(2\sigma_0^2)}\)，所以 \(\langle x\rangle=0\)、\(\Delta x=\sigma_0\)。Fourier变换给动量高斯谱，其标准差为 \(\Delta p=\hbar/(2\sigma_0)\)，初始态满足最小不确定度积 \(\Delta x\,\Delta p=\hbar/2\)。
(2) 自由Hamilton量 \(H=p^2/(2m)\) 使每个动量分量积累相位 \(e^{-ip^2t/(2m\hbar)}\)。完成高斯积分得 \(\psi(x,t)=(2\pi\sigma_0^2)^{-1/4}(1+i\tau)^{-1/2}\exp[-x^2/(4\sigma_0^2(1+i\tau))]\)，其中 \(\tau=\hbar t/(2m\sigma_0^2)\)。该式在 \(t=0\) 恢复初态。
(3) 波函数模平方仍是以原点为中心的高斯分布，其位置标准差 \(\sigma_x(t)=\sigma_0\sqrt{1+\tau^2}\)，故 \(\langle x\rangle=0\)、\((\Delta x)^2=\sigma_0^2+\hbar^2t^2/(4m^2\sigma_0^2)\)。自由Hamilton量与动量对易，动量分布保持不变，\((\Delta p)^2=\hbar^2/(4\sigma_0^2)\)。
(4) 当 \(t\to0\) 时位置方差恢复 \(\sigma_0^2\)；当 \(|t|\gg2m\sigma_0^2/\hbar\)，有 \(\Delta x\simeq\hbar|t|/(2m\sigma_0)\)，即自由波包线性展宽。由于不同动量分量传播速度不同，空间概率逐渐展开，但总动量以及动量分布仍守恒，不存在冲突。
